% -------------------------------------------------------------------------
% WBS ID: RES-MOD-EQS
% Deliverable: Governing Equations and State Variables
% Status: Draft complete
% Evidence status: Research specification complete for regional and local hydrodynamic equation families
% Review status: Not reviewed
% Last updated: 2026-07-18
% Dependencies:
% Downstream interfaces:
% -------------------------------------------------------------------------

\section{RES-MOD-EQS --- Governing Equations and State Variables}
\label{sec:res-mod-eqs}

\subsection{Purpose and Scope}
\label{sec:res-mod-eqs-purpose}

% TODO: Define what this Deliverable covers and explicitly excludes.

\subsection{Research Questions}
\label{sec:res-mod-eqs-research-questions}

% TODO: State the questions that must be answered before the Deliverable
% can be accepted.

\subsection{Terminology and Definitions}
\label{sec:res-mod-eqs-definitions}

% TODO: Define the technical terminology, symbols, quantities and
% classifications used in this Deliverable.

\subsection{Literature Evidence}
\label{sec:res-mod-eqs-literature-evidence}

% TODO: Record evidence from peer-reviewed papers, authoritative datasets,
% standards, technical reports and primary documentation.
%
% Do not create a detached literature-summary list. Explain how each source
% supports, contradicts or limits a project decision.

\subsubsection{Saito et al. (2014)}
\label{sec:res-mod-eqs-evidence-saito-2014}

\textcite{saito2014dispersion} compared nonlinear dispersive,
linear dispersive, and nonlinear non-dispersive models of the 2011
Tōhoku tsunami while retaining a common source, computational domain,
bathymetry, and observational dataset. The study showed that
frequency dispersion was required to reproduce short-period offshore
wave components generated by a spatially concentrated source, while
nonlinear advection and friction became increasingly important during
nearshore transformation and coastal reflection.

The paper is retained as the primary evidence that the nonlinear
shallow-water equations are necessary for nearshore propagation but
may be insufficient where the earthquake source retains physically
important short-wavelength components.

\subsubsection{Qin et al. (2018)}
\label{sec:res-mod-eqs-evidence-qin-2018}

\textcite{qin2018comparison} compared a depth-averaged nonlinear
shallow-water model with a three-dimensional Reynolds-averaged
Navier--Stokes model. The depth-averaged formulation reproduced several
bulk inundation quantities but could not represent the vertical
velocity variation and nonhydrostatic pressure associated with initial
impact. The three-dimensional formulation also permitted hydrodynamic
force to be obtained directly by integrating the fluid traction over
the structural surface rather than through an empirical drag relation.

The paper is retained as the primary evidence that the local barrier
domain requires a three-dimensional parent equation system when
pressure, traction, and structural loading are project outputs.

\subsection{Governing Relations and Quantitative Definitions}
\label{sec:res-mod-eqs-governing-relations}

% TODO: Record relevant equations, dimensionless groups, empirical models,
% extraction rules or quantitative definitions.
%
% Use align environments for displayed equations.
% Every displayed equation must have a unique \label{eq:...}.
% Do not place punctuation after displayed equations.

\subsection{System Boundary}
\label{sec:res-mod-eqs-system-boundary}

% TODO: Complete this RES-MOD Domain-specific subsection.

\subsection{State Variables and Parameters}
\label{sec:res-mod-eqs-state-variables-parameters}

% TODO: Complete this RES-MOD Domain-specific subsection.

\subsection{Continuous Governing Model}
\label{sec:res-mod-eqs-continuous-governing-model}

% TODO: Complete this RES-MOD Domain-specific subsection.

\subsubsection{Consolidated Governing Systems}
\label{sec:res-mod-eqs-consolidated-governing-systems}

The current research phase records three governing systems:

\begin{enumerate}
  \item a regional two-dimensional nonlinear shallow-water equation
    system;
  \item a local three-dimensional nonhydrostatic incompressible
    immiscible two-phase Navier--Stokes system;
  \item a transient structural-dynamics system with rigid, one-way
    and two-way FSI fidelity levels.
\end{enumerate}

The regional conserved state is:

\begin{align}
  \mathbf{U}_{\mathrm{r}}
  &=
  \left[
    h,\,
    q_x,\,
    q_y
  \right]^{\mathsf{T}}
  \label{eq:res-mod-eqs-regional-conserved-state}
  \\
  q_x
  &=
  h u
  \label{eq:res-mod-eqs-regional-x-discharge}
  \\
  q_y
  &=
  h v
  \label{eq:res-mod-eqs-regional-y-discharge}
\end{align}

where $h$ is water depth, $u$ and $v$ are depth-averaged horizontal
velocities, $q_x$ and $q_y$ are depth-integrated momenta, $b$ is bed
elevation positive upward, and $\eta=h+b$ is free-surface elevation.
The primitive regional variables are:

\begin{align}
  \mathbf{W}_{\mathrm{r}}
  &=
  \left[
    h,\,
    u,\,
    v,\,
    \eta,\,
    b
  \right]^{\mathsf{T}}
  \label{eq:res-mod-eqs-regional-primitive-state}
  \\
  h
  &=
  \eta
  -
  b
  \label{eq:res-mod-eqs-regional-depth-from-surface-bed}
\end{align}

The regional solve is horizontal and depth averaged in $(x,y,t)$.
Vertical sections are diagnostic visualisations extracted from this
two-dimensional state; they are not separate solved cross-sectional
models.

The selected regional baseline is the nonlinear shallow-water balance
law:

\begin{align}
  \frac{\partial \mathbf{U}_{\mathrm{r}}}{\partial t}
  +
  \frac{\partial \mathbf{F}_{\mathrm{r}}}{\partial x}
  +
  \frac{\partial \mathbf{G}_{\mathrm{r}}}{\partial y}
  &=
  \mathbf{S}_{\mathrm{b}}
  +
  \mathbf{S}_{\mathrm{f}}
  +
  \mathbf{S}_{\mathrm{sp}}
  \label{eq:res-mod-eqs-regional-nlswe-balance}
\end{align}

with fluxes:

\begin{align}
  \mathbf{F}_{\mathrm{r}}
  &=
  \left[
    q_x,\,
    \frac{q_x^2}{h}
    +
    \frac{1}{2}g h^2,\,
    \frac{q_x q_y}{h}
  \right]^{\mathsf{T}}
  \label{eq:res-mod-eqs-regional-x-flux}
  \\
  \mathbf{G}_{\mathrm{r}}
  &=
  \left[
    q_y,\,
    \frac{q_x q_y}{h},\,
    \frac{q_y^2}{h}
    +
    \frac{1}{2}g h^2
  \right]^{\mathsf{T}}
  \label{eq:res-mod-eqs-regional-y-flux}
\end{align}

The bed-slope source term is:

\begin{align}
  \mathbf{S}_{\mathrm{b}}
  &=
  \left[
    0,\,
    -g h \frac{\partial b}{\partial x},\,
    -g h \frac{\partial b}{\partial y}
  \right]^{\mathsf{T}}
  \label{eq:res-mod-eqs-regional-bed-slope-source}
\end{align}

The Manning-type bottom-friction term is retained as:

\begin{align}
  \mathbf{S}_{\mathrm{f}}
  &=
  \left[
    0,\,
    -g n_{\mathrm{M}}^2
    \frac{
      \sqrt{
        q_x^2
        +
        q_y^2
      }
      q_x
    }{
      h^{7/3}
    },\,
    -g n_{\mathrm{M}}^2
    \frac{
      \sqrt{
        q_x^2
        +
        q_y^2
      }
      q_y
    }{
      h^{7/3}
    }
  \right]^{\mathsf{T}}
  \label{eq:res-mod-eqs-regional-bottom-friction-source}
\end{align}

Open-boundary sponge damping is part of the established regional model
where artificial boundary reflections would otherwise contaminate the
domain:

\begin{align}
  \mathbf{S}_{\mathrm{sp},i}
  &=
  -\sigma_i
  \left[
    h_i-h_{\mathrm{ref},i},\,
    q_{x,i},\,
    q_{y,i}
  \right]^{\mathsf{T}}
  \label{eq:res-mod-eqs-regional-sponge-source}
  \\
  h_{\mathrm{ref},i}
  &=
  \max
  \left(
    \eta_{\mathrm{ref}}
    -
    b_i,\,
    0
  \right)
  \label{eq:res-mod-eqs-regional-sponge-reference-depth}
\end{align}

Earthquake forcing is represented by dynamic bathymetry:

\begin{align}
  b(x,y,t)
  &=
  b_0(x,y)
  +
  \sum_m
  \Delta b_m(x,y)R_m(t)
  \label{eq:res-mod-eqs-regional-dynamic-bed}
\end{align}

Because $h$ is the conserved depth variable in
\cref{eq:res-mod-eqs-regional-conserved-state}, this formulation does
not add a separate earthquake mass source to the depth equation.

The Coriolis source term remains an optional sensitivity term rather
than part of the baseline locked regional model:

\begin{align}
  \mathbf{S}_{\mathrm{c}}
  &=
  \left[
    0,\,
    f q_y,\,
    -f q_x
  \right]^{\mathsf{T}}
  \label{eq:res-mod-eqs-regional-coriolis-source}
\end{align}

where the sign convention follows the chosen horizontal coordinate
orientation.

Wetting and drying are part of the regional model requirements rather
than a separate governing-equation family. A dispersive extension
remains optional and shall be activated only where regional comparison
shows that nonlinear shallow-water propagation does not retain the
required phase, amplitude, or spectral content.

The local fluid governing system remains the three-dimensional
nonhydrostatic incompressible immiscible two-phase Navier--Stokes
system defined in
\cref{sec:res-mod-eqs-local-selected-equations}. Its required baseline
outputs include breaking, run-up, overtopping, separation,
recirculation, pressure and shear traction, URANS turbulence variables
and optional LES comparison fields. Moving fluid--structure boundaries
belong to deferred FSI activation and are not part of the fixed rigid
barrier baseline.

The structural governing system is transient structural dynamics. Its
fidelity hierarchy is rigid barrier, one-way CFD-to-FEA, and two-way
FSI where feedback is material. The structural formulation shall not be
permanently reduced to a single degree of freedom.

The reviewed Tōhoku hindcasts do not support the use of one governing
model across every spatial and physical regime. A provisional
model hierarchy is therefore defined as:

\begin{align}
  \mathcal{M}
  &=
  \begin{cases}
    \mathcal{M}_{\mathrm{source}},
    &
    \text{earthquake-driven tsunami generation}
    \\
    \mathcal{M}_{\mathrm{disp}},
    &
    \text{propagation where frequency dispersion is material}
    \\
    \mathcal{M}_{\mathrm{NLSW}},
    &
    \text{nonlinear nearshore propagation and inundation}
    \\
    \mathcal{M}_{\mathrm{3D}},
    &
    \text{local barrier interaction and vertically resolved flow}
  \end{cases}
  \label{eq:res-mod-eqs-provisional-model-hierarchy}
\end{align}

The hierarchy describes separable modelling regimes rather than a
requirement to implement four independent solvers. A sufficiently
general model may cover more than one regime, provided that its
accuracy and computational cost remain acceptable.

\subsubsection{Selected Three-Dimensional Fluid Equations}
\label{sec:res-mod-eqs-local-selected-equations}

The selected parent model is the three-dimensional, nonhydrostatic,
incompressible, immiscible, two-phase Navier--Stokes system in a
one-fluid representation.

Mass conservation is:

\begin{align}
  \nabla\cdot\mathbf{u}_{\mathrm{f}}
  &=
  0
  \label{eq:res-mod-eqs-local-continuity}
\end{align}

Momentum conservation is:

\begin{align}
  \frac{\partial
    \left(
      \rho_{\mathrm{f}}\mathbf{u}_{\mathrm{f}}
  \right)}
  {\partial t}
  +
  \nabla\cdot
  \left(
    \rho_{\mathrm{f}}
    \mathbf{u}_{\mathrm{f}}
    \otimes
    \mathbf{u}_{\mathrm{f}}
  \right)
  &=
  \nabla\cdot\boldsymbol{\sigma}_{\mathrm{f}}
  +
  \rho_{\mathrm{f}}\mathbf{g}
  +
  \mathbf{f}_{\sigma}
  +
  \mathbf{f}_{\mathrm{ext}}
  \label{eq:res-mod-eqs-local-momentum}
\end{align}

The fluid stress tensor is:

\begin{align}
  \boldsymbol{\sigma}_{\mathrm{f}}
  &=
  -p\mathbf{I}
  +
  2\mu_{\mathrm{eff}}
  \mathbf{D}
  \left(
    \mathbf{u}_{\mathrm{f}}
  \right)
  \label{eq:res-mod-eqs-local-fluid-stress}
  \\
  \mathbf{D}
  \left(
    \mathbf{u}_{\mathrm{f}}
  \right)
  &=
  \frac{1}{2}
  \left[
    \nabla\mathbf{u}_{\mathrm{f}}
    +
    \left(
      \nabla\mathbf{u}_{\mathrm{f}}
    \right)^{\mathsf{T}}
  \right]
  \label{eq:res-mod-eqs-local-rate-of-deformation}
\end{align}

The water phase is transported according to:

\begin{align}
  \frac{\partial\alpha}{\partial t}
  +
  \nabla\cdot
  \left(
    \alpha\mathbf{u}_{\mathrm{f}}
  \right)
  &=
  0
  \label{eq:res-mod-eqs-local-phase-transport}
\end{align}

The mixture density and molecular viscosity are:

\begin{align}
  \rho_{\mathrm{f}}
  &=
  \alpha\rho_{\mathrm{w}}
  +
  \left(
    1-\alpha
  \right)\rho_{\mathrm{a}}
  \label{eq:res-mod-eqs-local-mixture-density}
  \\
  \mu_{\mathrm{f}}
  &=
  \alpha\mu_{\mathrm{w}}
  +
  \left(
    1-\alpha
  \right)\mu_{\mathrm{a}}
  \label{eq:res-mod-eqs-local-mixture-viscosity}
\end{align}

The effective viscosity is expressed generically as:

\begin{align}
  \mu_{\mathrm{eff}}
  &=
  \mu_{\mathrm{f}}
  +
  \mu_{\mathrm{t}}
  \label{eq:res-mod-eqs-local-effective-viscosity}
\end{align}

where $\mu_{\mathrm{t}}$ is supplied by the selected turbulence
treatment.

The term $\mathbf{f}_{\sigma}$ represents surface-tension forcing.
Its physical contribution is expected to be small at tsunami and
barrier scales where the Weber number is large. It remains in the
general equation set so that its omission can be justified through a
regime assessment rather than imposed implicitly.

The term $\mathbf{f}_{\mathrm{ext}}$ permits later local extensions,
including locally activated porous resistance or other justified body
forces. It is zero in the initial impermeable-barrier model.

Decision trace: the equation set in
\cref{eq:res-mod-eqs-local-continuity,eq:res-mod-eqs-local-momentum,eq:res-mod-eqs-local-phase-transport}
is selected as the local parent hydrodynamic model. Source role:
`3D-RANSVOF`. Supporting sources:
\textcite{qin2018comparison} and \textcite{watanabe2022validation}.
Project-specific conclusion: a three-dimensional nonhydrostatic
air--water model is required when pressure, shear, force, moment,
run-up and overtopping are outputs. Limitation: the continuous equation
selection does not validate the selected mesh, VOF advection, pressure
correction, turbulence closure or wall treatment. Downstream handoff:
`RES-NUM-SPEC` shall define the finite-volume solution procedure and
`RES-VER-MET` shall define acceptance metrics for the extracted
hydrodynamic quantities.

\subsubsection{Required Hydrodynamic Outputs}
\label{sec:res-mod-eqs-local-required-outputs}

The local model shall preserve the spatial fluid traction:

\begin{align}
  \mathbf{t}_{\mathrm{f}}
  &=
  \boldsymbol{\sigma}_{\mathrm{f}}\mathbf{n}
  \label{eq:res-mod-eqs-local-fluid-traction}
\end{align}

The resultant force and moment acting on a structural surface
$\Gamma_{\mathrm{fs}}$ are:

\begin{align}
  \mathbf{F}_{\mathrm{f}}(t)
  &=
  \int_{\Gamma_{\mathrm{fs}}}
  \mathbf{t}_{\mathrm{f}}
  \,\mathrm{d}\Gamma
  \label{eq:res-mod-eqs-local-resultant-force}
  \\
  \mathbf{M}_{\mathrm{f}}(t)
  &=
  \int_{\Gamma_{\mathrm{fs}}}
  \left(
    \mathbf{x}-\mathbf{x}_{0}
  \right)
  \times
  \mathbf{t}_{\mathrm{f}}
  \,\mathrm{d}\Gamma
  \label{eq:res-mod-eqs-local-resultant-moment}
  \\
  \mathbf{J}_{\mathrm{f}}
  &=
  \int_{t_{0}}^{t_{1}}
  \mathbf{F}_{\mathrm{f}}(t)
  \,\mathrm{d}t
  \label{eq:res-mod-eqs-local-force-impulse}
\end{align}

The model shall output the complete pressure and traction fields rather
than only integrated force coefficients. This is required for
conservative transfer to a structural mesh and for calculating:

\begin{itemize}
  \item local peak pressure;
  \item pressure duration;
  \item resultant force;
  \item force impulse;
  \item overturning and torsional moments;
  \item centre-of-pressure migration;
  \item transmitted discharge and momentum;
  \item overtopping volume;
  \item reflected and transmitted wave energy;
  \item downstream velocity and turbulence fields.
\end{itemize}

\subsection{Local Three-Dimensional Governing-Equation Selection}
\label{sec:res-mod-eqs-local-three-dimensional-selection}

The regional depth-averaged model is intended to reproduce tsunami
generation, propagation, nearshore transformation, and the incident
state supplied to the local barrier domain. It is not intended to
resolve the complete local interaction between the wave, barrier,
free surface, and structural surface.

The evidence retained for this Deliverable establishes two distinct
conclusions:

\begin{enumerate}
  \item a resolved three-dimensional nonhydrostatic model is required
    for local barrier interaction, pressure, traction, and
    structural loading;
  \item selection of physically complete equations does not remove
    the requirement to validate the turbulence closure,
    free-surface treatment, boundary conditions, and numerical
    implementation.
\end{enumerate}

\subsubsection{Candidate Parent Equation Families}
\label{sec:res-mod-eqs-local-candidate-families}

The candidate formulations were ranked according to their ability to
represent:

\begin{itemize}
  \item three-dimensional and vertical fluid motion;
  \item nonhydrostatic pressure;
  \item wave breaking, run-up, overtopping, and separation;
  \item spatially and temporally resolved structural traction;
  \item moving fluid--structure interfaces;
  \item the current rigid, impermeable, and well-ventilated barrier
    design space;
  \item later multiphysics extensions without replacing the parent
    fluid model.
\end{itemize}

Published wall-clock times were not included in the equation-family
ranking. Computational performance depends on the numerical method,
linear and pressure solvers, spatial resolution, hardware, parallel
architecture, memory movement, and implementation quality. These
quantities will be assessed during numerical-method and software
selection rather than inferred from unrelated published simulations.

\begin{table}[htbp]
  \centering
  \small
  \renewcommand{\arraystretch}{1.25}
  \begin{tabular}{
      p{0.07\textwidth}
      p{0.27\textwidth}
      p{0.28\textwidth}
      p{0.28\textwidth}
    }
    \hline
    Rank
    &
    Formulation
    &
    Principal capability
    &
    Project decision
    \\
    \hline
    1
    &
    Three-dimensional incompressible, immiscible,
    two-phase Navier--Stokes
    &
    Resolves vertical acceleration, nonhydrostatic pressure,
    free-surface deformation, flow separation, and structural
    traction
    &
    Selected as the parent local fluid equation set
    \\
    2
    &
    Compressible two-phase Navier--Stokes
    &
    Adds gas and liquid compressibility, acoustic propagation,
    and trapped-gas compression
    &
    Deferred unless enclosed-air or compressibility effects become
    material
    \\
    3
    &
    Three-dimensional nonhydrostatic multilayer equations
    &
    Introduces vertical structure at lower dimensional complexity
    than fully resolved Navier--Stokes
    &
    Not selected for barrier-scale pressure and FSI owing to
    limited direct load-validation evidence
    \\
    4
    &
    Three-dimensional hydrostatic multilayer equations
    &
    Resolves several vertical layers while retaining hydrostatic
    pressure
    &
    Rejected for local impulsive and nonhydrostatic impact
    \\
    \hline
  \end{tabular}
  \caption{Project-specific ranking of candidate parent fluid
  equation families for the local barrier domain}
  \label{tab:res-mod-eqs-local-parent-ranking}
\end{table}

Reynolds-averaged Navier--Stokes, large-eddy simulation, and hybrid
RANS--LES are not treated as separate parent conservation-law systems.
They are alternative averaging, filtering, and closure treatments
attached to the selected Navier--Stokes equations.

Volume-averaged Reynolds-averaged Navier--Stokes is similarly treated
as a porous-region extension rather than as the default parent model.

Particle methods such as smoothed particle hydrodynamics and the moving
particle semi-implicit method are numerical formulations. They will be
evaluated during numerical-method selection and are not candidates in
the present governing-equation ranking.

\subsubsection{Selected Fluid State Variables}
\label{sec:res-mod-eqs-local-state-variables}

The operational local URANS state is defined as:

\begin{align}
  \mathcal{Q}_{\mathrm{f}}
  &=
  \left\{
    \mathbf{u}_{\mathrm{f}},
    p,
    \alpha,
    k,
    \omega
  \right\}
  \label{eq:res-mod-eqs-local-fluid-state}
\end{align}

where:

\begin{itemize}
  \item $\mathbf{u}_{\mathrm{f}}$ is the three-dimensional fluid
    velocity;
  \item $p$ is the mechanical pressure;
  \item $\alpha$ is the water volume fraction;
  \item $k$ is turbulent kinetic energy;
  \item $\omega$ is the specific dissipation rate.
\end{itemize}

The baseline closure attached to this state is the $k$--$\omega$ SST
URANS closure. Source role: `3D-SST`. Supporting source:
\textcite{menter1994twoequation}. Project-specific conclusion: the
high-Reynolds-number operational local model shall retain explicit
turbulence variables rather than treating the simulation as
molecular-viscosity-only flow. Limitation: LES and hybrid RANS--LES
remain higher-fidelity or later-comparison modes and require their own
state variables when activated. Downstream handoff: `RES-MOD-SRC` owns
the turbulence closure equations and `RES-NUM-SPEC` owns their discrete
solution order.

\subsection{Source Terms and Closures}
\label{sec:res-mod-eqs-source-terms-closures}

% TODO: Complete this RES-MOD Domain-specific subsection.

\subsection{Initial, Boundary and Interface Conditions}
\label{sec:res-mod-eqs-initial-boundary-interface-conditions}

% TODO: Complete this RES-MOD Domain-specific subsection.

\subsection{Model Coupling Requirements}
\label{sec:res-mod-eqs-model-coupling-requirements}

% TODO: Complete this RES-MOD Domain-specific subsection.

\subsection{Sufficient Conditions for Validity}
\label{sec:res-mod-eqs-sufficient-conditions-validity}

% TODO: Complete this RES-MOD Domain-specific subsection.

\subsection{Candidate Approaches}
\label{sec:res-mod-eqs-candidate-approaches}

% TODO: Define the credible candidate approaches considered.

\subsection{Critical Comparison}
\label{sec:res-mod-eqs-critical-comparison}

% TODO: Compare candidates using physically and project-relevant criteria.
% Do not select an approach solely on popularity or implementation ease.

\subsection{Assumptions and Applicability}
\label{sec:res-mod-eqs-assumptions}

% TODO: State assumptions and the conditions under which they are valid.

\subsection{Limitations and Failure Conditions}
\label{sec:res-mod-eqs-limitations}

% TODO: State where the evidence, model or selected approach ceases to be
% reliable or applicable.

\subsection{Project-Specific Conclusions}
\label{sec:res-mod-eqs-conclusions}

% TODO: Convert the evidence into conclusions specific to the Studentship
% project. Do not merely restate literature findings.

The current review selects the parent local equation family and defines
the regional incident-wave model that will provide the state supplied
to that model.

The provisional regional baseline is a nonlinear depth-averaged model
with optional dispersive terms. Higher-order far-field corrections
shall remain independent modules and shall not be included in the
near-Japan benchmark unless observational comparison establishes that
they materially affect the required outputs.

\subsubsection{Project-Specific Governing-Equation Decision}
\label{sec:res-mod-eqs-local-project-decision}

The initial hydrodynamic architecture shall contain two principal
equation families:

\begin{align}
  \mathcal{M}_{\mathrm{hydro}}
  &=
  \left\{
    \mathcal{M}_{\mathrm{regional}},
    \mathcal{M}_{\mathrm{local}}
  \right\}
  \label{eq:res-mod-eqs-two-model-architecture}
  \\
  \mathcal{M}_{\mathrm{regional}}
  &:=
  \text{selected depth-averaged regional formulation}
  \label{eq:res-mod-eqs-regional-model-definition}
  \\
  \mathcal{M}_{\mathrm{local}}
  &:=
  \text{three-dimensional incompressible two-phase
  Navier--Stokes formulation}
  \label{eq:res-mod-eqs-local-model-definition}
\end{align}

The local model shall initially represent a rigid, impermeable, and
well-ventilated barrier. Turbulence closure, FSI coupling, porous
resistance, and compressible-air treatment are modular attachments
whose activation is governed by the physical regime and design
geometry.

For the baseline hydrodynamic campaign the closure attachment is no
longer optional: the operational local model is URANS with
$k$--$\omega$ SST, while LES is retained for selected high-fidelity
comparison cases.

\subsection{Model Selection Criteria}
\label{sec:res-mod-eqs-criteria}

The regional governing model shall be selected using:

\begin{itemize}
  \item the spatial and temporal bandwidth of the earthquake source;
  \item propagation distance;
  \item local water depth;
  \item expected wave nonlinearity;
  \item the importance of wetting and drying;
  \item the required accuracy of arrival time, amplitude, phase, and
    spectral content;
  \item the intended coupling location with the local
    three-dimensional model;
  \item the available computational resources.
\end{itemize}

The baseline comparison shall evaluate:

\begin{align}
  \mathcal{M}_{\mathrm{NLSW}}
  \qquad\text{and}\qquad
  \mathcal{M}_{\mathrm{NLSW+B}}
  \label{eq:res-mod-eqs-baseline-comparison}
\end{align}

where $\mathcal{M}_{\mathrm{NLSW+B}}$ denotes a nonlinear
shallow-water formulation augmented by the selected Boussinesq or
equivalent dispersive terms.

The dispersive extension shall be retained only where it produces a
material improvement in the benchmark quantities relative to its
additional computational cost.

\subsection{Selected Approach and Rationale}
\label{sec:res-mod-eqs-selected-approach}

% TODO: Record the selected approach only after sufficient evidence exists.
% State the alternatives rejected and the reasons for rejection.

\subsection{Unresolved Decisions}
\label{sec:res-mod-eqs-unresolved-decisions}

% TODO: Record unresolved questions, required evidence and decision owners.

\subsection{Requirements and Handoffs}
\label{sec:res-mod-eqs-requirements}

% TODO: State requirements passed to other Research Domains, Software,
% verification activities or later execution work.

The following conclusions are retained from the comparative studies:

\begin{enumerate}
  \item nonlinear momentum transport and friction are required in
    the nearshore region and for waves that have interacted with
    the coastline;
  \item frequency dispersion becomes material when the source
    contains short-wavelength components or when small
    phase-speed differences accumulate over sufficient distance;
  \item a source model that removes short-wavelength content may
    conceal the importance of dispersion;
  \item time-dependent source generation may alter arrival time and
    spectral content relative to instantaneous initialisation;
  \item seawater-density stratification, elastic loading, and
    gravitational-potential change principally improve
    trans-oceanic phase and arrival-time accuracy;
  \item no far-field correction shall be enabled solely on the basis
    that it increases physical complexity.
\end{enumerate}

\subsection{Deliverable Completion Record}
\label{sec:res-mod-eqs-completion-record}

% TODO: Record whether the Deliverable is:
%
% - Not started
% - In progress
% - Draft complete
% - Reviewed
% - Accepted
%
% Acceptance requires:
%
% - the research questions to be answered;
% - claims to be supported by appropriate evidence;
% - assumptions and limitations to be explicit;
% - the project-specific conclusion to be stated;
% - downstream requirements to be traceable;
% - unresolved decisions to be closed or formally deferred.
